% !TeX program = xelatex
% UTF-8 transcription of the three photographs in images/pic01.jpg--pic03.jpg.
% Compile directly with XeLaTeX, or input from a book using XePersian.
% Yellow emphasis follows the handwriting; uncertain readings are marked.
% The source folder contains three photographs, so there is no pic04.jpg.
% pic01.jpg = 20261010_203440294.jpg (handwritten page 1)
% pic02.jpg = 20261010_203444703.jpg (handwritten page 2)
% pic03.jpg = 20261010_203448179.jpg (handwritten page 3)
\newif\ifMNDStandalone
\ifdefined\settextfont
  \MNDStandalonefalse
\else
  \MNDStandalonetrue
\fi
\ifMNDStandalone
  \documentclass[a4paper,12pt]{extbook}
  \usepackage[margin=2.5cm]{geometry}
  \usepackage{xcolor}
  \usepackage{amsmath}
  \usepackage{enumitem}
  \usepackage{xepersian}
  \settextfont{Amiri}
  \setlatintextfont{FreeSerif}
  \pagestyle{plain}
\fi
\ifMNDStandalone
  \begin{document}
\fi
\begingroup
\emergencystretch=2em
\newcommand{\MNDen}[1]{\lr{#1}}
\ifdefined\colorbox
  \newcommand{\MNDhl}[1]{\colorbox{yellow!35}{\strut #1}}
\else
  \newcommand{\MNDhl}[1]{\textbf{#1}}
\fi
\newcommand{\MNDunclear}[1]{[خوانش نامطمئن: #1]}
\newcommand{\MNDpage}[2]{\clearpage\section*{صفحهٔ #1 ـ ۱۸: بیماری‌های نورون حرکتی}\noindent{\footnotesize\MNDen{#2}}\par\medskip}

\chapter{بیماری‌های نورون حرکتی (\MNDen{MND})}
\noindent
رونویسی از سه صفحهٔ دست‌نویس فارسی و انگلیسی. هایلایت‌ها از اصل یادداشت حفظ شده‌اند؛ خوانش‌های نامطمئن مشخص شده‌اند.
% Printed stationery on each image: Date: / / ; Subject: ;
% Sat. Sun. Mon. Tue. Thu. Wed. Fri. (in the order printed).
% The date and weekday boxes are blank; the handwritten subject is preserved below.

\MNDpage{۱}{pic01.jpg — 20261010\_203440294.jpg}
\section*{بیماری‌های نورون حرکتی}
\begin{itemize}
\item به بیماری‌هایی گفته می‌شود که در آن‌ها نورون‌های حرکتی \MNDen{CNS} از بین رفته‌اند.
\item بیماری‌های نورون حرکتی تنها به نورون حرکتی در نواحی زیر مربوط‌اند:
\begin{enumerate}
\item شاخ قدامی مادهٔ خاکستری نخاع: نورون، از شاخ قدامی نخاع تا عضلات ارادی اندام‌ها.
\item ساقهٔ مغز: این نورون‌ها، هسته‌های حرکتی بعضی اعصاب کرانیال هستند.
\end{enumerate}
\item نورون محرکهٔ فوقانی: نورون‌های حرکتی کورتکس فرونتال. این نورون‌ها در جلوی شیار رولاندیک قرار داشته و آکسون آن‌ها به نورون‌های حرکتی ساقهٔ مغز و نخاع ختم می‌شود.
\item علائم ضایعات \MNDen{LMN}: \MNDhl{ضعف و فلجی} / آتروفی شدید / کاهش یا از بین رفتن رفلکس تاندون / فاسیکولاسیون.
\item علائم ضایعات \MNDen{UMN}: ضعف / \MNDhl{اسپاستیته} / \MNDhl{افزایش \MNDen{DTR}} / \MNDhl{بابنسکی} / عدم وجود آتروفی.
\item در بیماری‌های \MNDen{MN}، \MNDhl{اختلال اسفنکتر}، \MNDhl{دمانس}، \MNDhl{اختلال حسی}، \MNDhl{مخچه‌ای} و \MNDhl{اکستراپیرامیدال} \MNDhl{دیده نمی‌شود}.
\item اتیولوژی بیماری‌های \MNDen{MN}: ارثی، عفونی، کمبودهای آنزیمی، بدخیمی.
\end{itemize}

\section*{انواع بیماری‌های نورون حرکتی}
\begin{itemize}
\item \MNDen{SMA} (اسپاینال ماسکولار آتروفی): درگیری فقط در نورون‌های حرکتی شاخ قدامی نخاع؛ \MNDhl{درگیری \MNDen{LMN}}.
\item \MNDen{PBP} (پروگرسیو بولبار پالزی): ضعف، آتروفی و فاسیکولاسیون در عضلات صورت، زبان، حلق و حنجره.
\begin{itemize}
\item در واقع عضلاتی که از هسته‌های حرکتی اعصاب کرانیال منشأ عصب می‌گیرند.
\end{itemize}
\item \MNDen{ALS}: \MNDhl{هم درگیری \MNDen{LMN}} و \MNDhl{هم \MNDen{UMN}}.
\end{itemize}

\MNDpage{۲}{pic02.jpg — 20261010\_203444703.jpg}
\section*{\MNDen{ALS}}
\begin{itemize}
\item \MNDhl{شایع‌ترین بیماری \MNDen{MN}}.
\item انواع:
\begin{itemize}
\item اسپورادیک: \MNDhl{شایع‌ترین نوع} / \MNDhl{\MNDen{$F<M$}} / \MNDhl{دههٔ پنجم تا هفتم}.
\item ارثی: \MNDhl{\MNDen{AD}} / سن پایین‌تر.
\begin{itemize}
\item \MNDen{$\tfrac{1}{2}$} موارد: ژن \MNDhl{آنزیم سوپراکسید دیس‌موتاز ۱}، کروموزوم ۲۱.
% The fraction is transcribed as written; this is not a medical correction.
\end{itemize}
\end{itemize}
\item اتیولوژی: نامشخص / سموم مثل \MNDunclear{سلنیوم} و جیوه / علل ایمونولوژیک / ژنتیک / ویروس‌ها.
\item پاتولوژی: تغییرات دژنراتیو در نورون‌های حرکتی نخاع، \MNDen{brainstem}، کورتکس حرکتی.
\item علائم: یک اختلال مختلط \MNDen{UMN + LMN}.
\begin{itemize}
\item اسپاستیته، \MNDhl{\MNDen{$DTR\uparrow$}}، \MNDhl{بابنسکی} / ضعف، \MNDhl{آتروفی}، فاسیکولاسیون.
\end{itemize}
\item \MNDhl{عضلات حرکتی چشم هرگز در \MNDen{ALS}} \MNDhl{درگیر نمی‌شوند}.
\item \MNDhl{بیماران علائم حسی ندارند}.
\item سیر بیماری: تدریجی و موذیانه / ضعف یک یا چند عضله / درگیری غیرقرینهٔ دیستال اندام‌ها.
\begin{itemize}
\item بیمار تدریجاً در استفاده از کلید، بستن دکمهٔ لباس، بلند کردن اجسام دچار مشکل می‌شود.
\end{itemize}
\item درگیری عضلات بولبار: اشکال در بلع، تو دماغی صحبت کردن، اختلال در سخن گفتن، ضعف و لاغری زبان، فاسیکولاسیون زبان.
\item دیسفاژی و اختلال تنفسی می‌توانند در اوایل بیماری هم بروز کنند.
\item علائم سودوبولبار در برخی بیماران: \MNDhl{گریه و خندهٔ پاتولوژیک}.
\item اسپاسم و کرامپ عضلانی، به‌خصوص در شب‌ها شایع‌تر است.
\item \MNDhl{مثانهٔ نوروژنیک در مراحل دیررس بیماری رخ می‌دهد}.
\item \MNDen{$\tfrac{1}{2}$} بیماران دچار درجاتی از تغییرات شناختی، رفتاری و شخصیتی می‌شوند.
\item دیسترس تنفسی در صورت درگیری عضلات تنفسی: نیاز به حمایت تنفسی.
\item تشخیص: فیبریلاسیون و فاسیکولاسیون در \MNDen{EMG}.
\item درمان: \MNDhl{ریلوزول}.
\end{itemize}

\MNDpage{۳}{pic03.jpg — 20261010\_203448179.jpg}
\section*{ادامهٔ \MNDen{ALS}}
\begin{itemize}
\item طول عمر بیماران بعد از تشخیص \MNDhl{۳–۴ سال} است.
\item شایع‌ترین علل مرگ: نارسایی تنفسی، عفونت‌های تنفسی یا سیستمیک.
\item پروگنوز بد: \MNDhl{سطح سرمی \MNDunclear{لیپید} \MNDen{$\downarrow$}}، کاهش شدید وزن، سن بالا، ظرفیت‌های ریوی \MNDen{$\downarrow$}.
\item یافته‌های کاراکتریستیک \MNDen{ALS}:
\begin{itemize}
\item فاسیکولاسیون زبان، \MNDhl{\MNDen{$DTR\uparrow$}}، آتروفی و تحلیل عضلات، دیسفاژی و دیزآرتری.
\item بابنسکی، آتروفی عضلات اینترنسیک.
\end{itemize}
\item علائمی که در \MNDen{ALS} وجود ندارد:
\begin{itemize}
\item دمانس، اختلال اسفنکتر، اختلال حسی، مخچه‌ای، اکستراپیرامیدال.
\end{itemize}
\end{itemize}

\section*{پولیومیلیت}
\begin{itemize}
\item نورون‌های حرکتی مادهٔ خاکستری شاخ قدامی نخاع / \MNDhl{\MNDen{LMN}}.
\item \MNDhl{فلج شل غیرقرینه}.
\item علائم: مونوپارزی، پاراپارزی یا کوادری‌پارزی شل.
\item فقدان یا \MNDen{$\downarrow DTR$}.
\item فقدان علائم حسی و اسفنکتری.
\end{itemize}

\section*{سندرم \MNDen{Post Polio}}
\begin{itemize}
\item گاهی چند سال پس از ابتلا به پولیومیلیت، با افزایش سن سلول‌های حرکتی باقی‌مانده در شاخ قدامی از بین می‌روند؛ اختلال حرکتی افزایش می‌یابد.
\end{itemize}

\endgroup
\ifMNDStandalone
  \expandafter\enddocument
\fi
